\documentclass[11pt,a4paper]{article}
\usepackage[T1]{fontenc}
\usepackage{lmodern}
\usepackage[margin=28mm]{geometry}
\usepackage{amsmath,amssymb,amsthm,mathtools}
\usepackage{microtype}
\usepackage{xcolor}
\usepackage[colorlinks=true,linkcolor=blue!55!black,citecolor=blue!55!black,urlcolor=blue!55!black]{hyperref}
\numberwithin{equation}{section}
\newtheorem{theorem}{Theorem}[section]
\newtheorem{lemma}[theorem]{Lemma}
\newtheorem{proposition}[theorem]{Proposition}
\newtheorem{corollary}[theorem]{Corollary}
\theoremstyle{remark}
\newtheorem{remark}[theorem]{Remark}
\newcommand{\R}{\mathbb R}
\newcommand{\T}{\mathbb T}
\newcommand{\Sph}{\mathbb S}
\newcommand{\la}[1]{\langle #1\rangle}
\newcommand{\calL}{\mathcal L}
\newcommand{\dd}{\,d}
\newcommand{\norm}[2]{\lVert #1\rVert_{#2}}
\DeclareMathOperator{\supp}{supp}
\title{Loss of rapid velocity decay in a broad smooth class\\
for the cutoff Boltzmann equation}
\author{Kaiwen Jin\\
\small\href{mailto:yui.ui.siki@gmail.com}{yui.ui.siki@gmail.com}}
\hypersetup{
 pdftitle={Loss of rapid velocity decay in a broad smooth class for the cutoff Boltzmann equation},
 pdfauthor={Kaiwen Jin}}
\date{15 October 2026}
\begin{document}
\maketitle
\begin{abstract}
We construct a spatially homogeneous, radial, nonnegative smooth solution of
the full-sphere hard-sphere Boltzmann equation whose initial value is a
Gaussian but whose velocity tail is a positive multiple of \(r^{-6}\) at
every positive time. Mass is one, and energy and entropy have uniform finite
upper bounds on a short time interval. This gives a negative answer to a
literal rapid-decay continuation statement when smoothness means joint
pointwise smoothness and no energy admissibility or positive-time Schwartz
condition is part of the solution class. The solution itself continues
smoothly in the critical weighted class: the example concerns loss of rapid
velocity decay and does not establish local derivative blowup or resolve
conservative continuation. The proof separates an explicit radial linear
part from a nonlinear remainder with one extra power of decay, and proves
nonnegativity by a compact-maximum argument. No novelty claim is made.
\end{abstract}

\section{Statement and interpretation}\label{sec:statement}

Consider
\begin{equation}\label{eq:boltzmann}
 \partial_t f+v\cdot\nabla_xf=Q(f,f),\qquad
 Q(F,F)(v)=\int_{\R^3}\int_{\Sph^2}|v-v_*|
 [F(v')F(v_*')-F(v)F(v_*)]\dd\sigma\dd v_*,
\end{equation}
where
\[
 v'=\frac{v+v_*}{2}+\frac{|v-v_*|}{2}\sigma,\qquad
 v_*'=\frac{v+v_*}{2}-\frac{|v-v_*|}{2}\sigma.
\]
The sphere has its usual surface measure of total area \(4\pi\);
\(\T^3=\R^3/\mathbb Z^3\). All collision factors in
\eqref{eq:boltzmann} have the same \(t,x\).

The formulation of AIM problem 366 asks whether a nonnegative smooth
solution on \([0,T)\times\T^3\times\R^3\), initially rapidly decreasing
together with all derivatives, extends beyond \(T\) with the same rapid
velocity decay if, uniformly in \(t,x\),
\begin{equation}\label{eq:macroscopic}
 m_0\leq\int_{\R^3}f\dd v\leq M_0,\qquad
 \int_{\R^3}|v|^2f\dd v\leq E_0,\qquad
 \int_{\R^3}f\log f\dd v\leq H_0.
\end{equation}
Here \(m_0,M_0,E_0>0\), \(H_0\in\R\), and \(0\log0=0\).
We address the literal reading in which smoothness is joint pointwise
\(C^\infty\) regularity. In particular, no energy inequality or locally
uniform positive-time Schwartz membership is assumed. The distinction
between these solution classes is essential to the conclusion.

For the time regularity assertion, let \(X_6\) be the space of continuous
radial profiles \(F:[0,\infty)\to\R\) with
\[
 \norm{F}{6}=\sup_{r\geq0}(1+r^2)^3|F(r)|<\infty.
\]
A homogeneous radial function is identified with its profile in this space.

\begin{theorem}\label{thm:construction}
There exist \(\varepsilon>0\), \(t_*>0\), and a nonnegative spatially
homogeneous radial solution \(f\) of \eqref{eq:boltzmann}, jointly
\(C^\infty\) on \([0,t_*]\times\T^3\times\R^3\), such that
\begin{gather}
 f(0,x,v)=\pi^{-3/2}e^{-|v|^2},\qquad
 \int_{\R^3}f(t,x,v)\dd v=1,\label{eq:datum-mass}\\
 \sup_{t,x}\int_{\R^3}|v|^2f(t,x,v)\dd v<\infty,\qquad
 \sup_{t,x}\int_{\R^3}f(t,x,v)\log f(t,x,v)\dd v<\infty,\label{eq:bounds}\\
 \lim_{r\to\infty}r^6f(t,x,r\omega)
       =4\varepsilon e^{-1/t^2}>0
 \quad(0<t\leq t_*,\ x\in\T^3,\ \omega\in\Sph^2).\label{eq:main-tail}
\end{gather}
All spatial and velocity derivatives of the initial datum are Schwartz.
The solution is \(C^\infty\) in time with values in \(X_6\).
\end{theorem}

Smoothness at a closed time endpoint means smooth one-sided traces.
The final \(t_*\) will lie strictly inside the interval on which the
solution is constructed. Thus the example already has a smooth
continuation in \(X_6\). For every \(0<T<t_*\), its restriction to
\([0,T)\) satisfies \eqref{eq:macroscopic}, but its endpoint at \(T\)
cannot be rapidly decreasing. No conclusion about blowup of local
derivatives, or about the conservative or energy-admissible continuation
problem, follows from this example.

The cutoff regularity question is discussed in
Silvestre's survey \cite[Section 8.5]{silvestre}. The formulation used
here is the explicit hard-sphere statement of AIM problem 366.\footnote{We
use this \href{https://github.com/MColbrook/AIM/blob/a153f469ee0ea64e876b2ad97c79d252c19a3e22/problems/366-cutoff-boltzmann-continuation.md}{archived version of the problem statement}.
The accompanying target comparison records its correspondence with the
note. Its interpretation, rather than a
claim to resolve every conventional formulation of conditional
regularity, is the point of the example.}
We use the account in Basile et al. \cite[Sections 1 and 7]{basile}
for the energy-increasing weak-solution precedent and its attribution
to Lu and Wennberg \cite{lu-wennberg}. Basile et al. also give a weak
construction. Our reference to Lu--Wennberg follows that account; their
full theorem statements remain unverified here. The account provides
context for contributions from arbitrarily large velocities. We do not
use a kinetic existence theorem from these works, and make no novelty
or priority claim.

The proof uses the leading radial operator
\[
 C_0F(r)=\frac4r\int_r^\infty sF(s)\dd s-rF(r).
\]
Its \(r^{-6}\) tail is critical: the two displayed terms cancel on that
power. After a cutoff and a mass correction, the difference between the
nonlinear collision operator and this leading part decays one power
faster. We prescribe a flat positive tail, solve for a faster-decaying
correction, and finally prove nonnegativity. Sections
\ref{sec:radial}--\ref{sec:positive} give these steps.

\section{Radial identities and leading terms}\label{sec:radial}

Write \(r=|v|\), \(e=r^2\), \(F(r)=F(v)\), and
\(G(e)=F(\sqrt e)\). We allow signed radial profiles until positivity has
been proved. Define
\[
 m(F)=4\pi\int_0^\infty r^2F(r)\dd r,\qquad
 \nu_F(v)=4\pi\int_{\R^3}|v-w|F(w)\dd w,
\]
and let \(Q^+(F_1,F_2)\) denote the gain integral with post-collision
factors \(F_1(v')F_2(v_*')\). It is a symmetric bilinear map whenever the integral is
absolute; symmetry follows by replacing \(\sigma\) with \(-\sigma\).
Then \(Q(F,F)=Q^+(F,F)-\nu_FF\).

For \(s>0\), put \(g_s(r)=\la r^{-s}\), where
\(\la r=(1+r^2)^{1/2}\), and set
\[
 X_s=\{F\in C([0,\infty)):\norm{F}{s}
       :=\sup_{r\geq0}\la r^s|F(r)|<\infty\}.
\]
Let \(X_s^0\) be the closed subspace on which
\(\la r^sF(r)\to0\). Multiplication by \(\la r^s\) identifies \(X_s\)
isometrically with the bounded continuous functions, so these spaces are
Banach spaces. Compactly supported continuous profiles are dense in
\(X_s^0\). Fix throughout
\begin{equation}\label{eq:p}
 p=\frac{13}{2},\qquad 6<p<7,\qquad
 M(r)=\pi^{-3/2}e^{-r^2}.
\end{equation}
The Gaussian has mass one and satisfies \(Q(M,M)=0\) by conservation of
the sum of the two squared speeds in a single collision. This identity
does not assert an energy balance for a time-dependent solution.
For \(j=0,1,2\),
\[
 \int_{\R^3}|v|^j|F(v)|\dd v
 \leq4\pi\norm{F}{6}\int_0^\infty r^{j+2}(1+r^2)^{-3}\dd r<\infty.
\]
Thus mass and the first two absolute moments are bounded on \(X_6\).
So are \(F\mapsto\int G(e)\dd e\) and
\(F\mapsto\int eG(e)\dd e\).

\subsection*{The exact collision formulas}
For \(e>0\), define the tail and local integrals
\begin{align*}
 A_e&=\int_e^\infty G(s)\dd s,&
 B_e&=\int_0^e\sqrt s\,G(s)\dd s,\\
 D_e&=\int_e^\infty sG(s)\dd s,&
 E_e&=\int_0^e s^{3/2}G(s)\dd s,\\
 C_e&=\int_0^e\dd s\int_{e-s}^e\dd t\,
        \sqrt{s+t-e}\,G(s)G(t).
\end{align*}
For every radial \(F\in X_6\),
\begin{align}
 \frac{Q^+(F,F)(r)}{8\pi^2}
   &=A_e^2+\frac{2B_eA_e}{r}+\frac{C_e}{r},\label{eq:gain}\\
 \frac{\nu_F(r)}{8\pi^2}
   &=rB_e+\frac{E_e}{3r}+D_e+\frac e3A_e.\label{eq:loss}
\end{align}
Here and below the formulas at \(r=0\) mean their continuous limits.

To check the normalization, the unsplit radial formula is
\begin{equation}\label{eq:min}
 \frac{Q(F,F)(r)}{8\pi^2}
 =\frac1r\int_0^\infty\dd e_*\int_0^{e+e_*}\dd e_1\,
 \min(\sqrt e,\sqrt{e_*},\sqrt{e_1},\sqrt{e_2})
 [G(e_1)G(e_2)-G(e)G(e_*)],
\end{equation}
where \(e_2=e+e_*-e_1\).
Indeed, for \(s=|v_*|\), let \(z\) be the cosine between \(v\) and
\(v_*\), \(U=|v+v_*|\), and \(d=|v-v_*|\).
If \(\theta\) is the angle between \(v+v_*\) and \(\sigma\), then
\[
 e_1=\frac{r^2+s^2+Ud\cos\theta}{2}.
\]
Thus \(e_1\) is uniform on an interval of length \(Ud\).
The two azimuthal integrations
give \(8\pi^2\,\dd z\,\dd e_1/(Ud)\). Multiplication by the hard-sphere
factor \(d\) and use of \(\dd z=U\,\dd U/(rs)\) leave
\(8\pi^2\,\dd U\,\dd e_1/(rs)\). For fixed output speeds the allowed
\(U\)-interval is
\[
 [|r-s|,r+s]\cap
 [|\sqrt{e_1}-\sqrt{e_2}|,\sqrt{e_1}+\sqrt{e_2}],
\]
of length \(2\min(r,s,\sqrt{e_1},\sqrt{e_2})\).
Finally \(s\,\dd s=\dd e_*/2\) gives \eqref{eq:min}.

For the gain, change variables from \((e_*,e_1)\) to
\((\xi,\zeta)=(e_1,e_2)\). The Jacobian is one,
\(e_*=\xi+\zeta-e\), and the domain is
\(\xi,\zeta\geq0\), \(\xi+\zeta\geq e\).
When both output energies exceed \(e\), the minimum kernel is \(\sqrt e=r\);
after division by \(r\) this region contributes \(A_e^2\).
If \(\xi\leq e\leq\zeta\), the minimum is \(\sqrt\xi\), and the
contribution is \(B_eA_e/r\). The exchanged region contributes the
same amount. In the remaining triangle
\(\xi,\zeta\leq e\), \(\xi+\zeta\geq e\), the minimum is
\(\sqrt{\xi+\zeta-e}\), giving \(C_e/r\).
These three contributions prove \eqref{eq:gain}.

For the loss, direct integration of the same kernel gives
\[
 \int_0^{e+e_*}
 \min(\sqrt e,\sqrt{e_*},\sqrt y,\sqrt{e+e_*-y})\dd y
 =\sqrt{\min(e,e_*)}
   \left[\max(e,e_*)+\frac{\min(e,e_*)}{3}\right].
\]
Indeed, split \(y\) at \(\min(e,e_*)\) and \(\max(e,e_*)\).
The two end integrals give \(4\min(e,e_*)^{3/2}/3\); the middle one
is \(\sqrt{\min(e,e_*)}\,|e-e_*|\).
After division by \(r\), the regions \(e_*\leq e\) and \(e_*\geq e\)
give respectively \(rB_e+E_e/(3r)\) and \(D_e+eA_e/3\).
This proves \eqref{eq:loss}.
All coordinate changes can first be made for nonnegative profiles.
The estimates below give finite majorants for the \(g_6\) envelope,
and hence justify the signed formulas by absolute integration.

There is no singularity at zero. Directly in the original operator,
\[
 Q^+(F,F)(0)=8\pi^2\left(\int_0^\infty G(s)\dd s\right)^2,\qquad
 \nu_F(0)=8\pi^2\int_0^\infty sG(s)\dd s.
\]
For later smoothness arguments, the local terms have the useful forms
\begin{align}
 \frac{B_e}{r}
   &=e\int_0^1\sqrt b\,G(eb)\dd b,&
 rB_e&=e^2\int_0^1\sqrt b\,G(eb)\dd b,\nonumber\\
 \frac{E_e}{3r}
   &=\frac{e^2}{3}\int_0^1b^{3/2}G(eb)\dd b,&
 \frac{C_e}{r}
   &=e^2\int_0^1\dd b\int_{1-b}^1\dd d\,
       \sqrt{b+d-1}\,G(eb)G(ed).\label{eq:origin}
\end{align}
They vanish to the stated orders at zero.

\subsection*{The mass correction and the faster remainder}
Choose a fixed smooth \(0\leq\chi\leq1\), zero on \([0,4]\) and one on
\([8,\infty)\). For example, with \(\eta(u)=e^{-1/u}\) for \(u>0\) and
zero otherwise, take
\(\chi(r)=\eta(r-4)/[\eta(r-4)+\eta(8-r)]\).
Recall that \(M\) is the unit-mass Gaussian in \eqref{eq:p}.
For \(r>0\), the leading expression is
\[
 C_0F(r)=\frac4r\int_r^\infty sF(s)\dd s-rF(r).
\]
Define
\begin{equation}\label{eq:cold}
 C_{\rm raw}F=\chi C_0F,\qquad
 \ell(F)=m(C_{\rm raw}F),\qquad
 CF=C_{\rm raw}F-\ell(F)M,\qquad \calL F=4\pi CF.
\end{equation}
The raw expression is zero near zero. For \(s>4\), \(\ell\) is bounded
on \(X_s\): the loss mass is controlled by
\(\int_4^\infty r^3|F(r)|\dd r\), and Tonelli gives the same bound for
the gain mass. In particular \(m(CF)=0\). The expression \(\calL F\)
is defined pointwise on \(X_s\); it need not belong to \(X_s\).

\begin{lemma}\label{lem:remainder}
The map
\[
 R(F)=Q(F,F)-4\pi m(F)CF
\]
is a continuous quadratic map \(X_6\to X_7\), with
\begin{align}
 \norm{R(F)}{7}&\leq C\norm{F}{6}^2,\nonumber\\
 \norm{R(F)-R(\widetilde F)}{7}
 &\leq C(\norm{F}{6}+\norm{\widetilde F}{6})
                \norm{F-\widetilde F}{6}.\label{eq:R}
\end{align}
It has zero mass. Since \(p<7\), its values belong to \(X_p^0\).
On the affine space \(m(F)=1\),
\begin{equation}\label{eq:split}
 Q(F,F)=\calL F+R(F).
\end{equation}
\end{lemma}
\begin{proof}
Put \(B_\infty=\int_0^\infty\sqrt s\,G(s)\dd s=m(F)/(2\pi)\).
For \(r\geq8\), subtraction of the uncorrected leading radial part from
\eqref{eq:gain}--\eqref{eq:loss}, after division by \(8\pi^2\), leaves
\begin{equation}\label{eq:remainder}
 A_e^2-\frac{2(B_\infty-B_e)A_e}{r}+\frac{C_e}{r}
 +F(r)\left[r(B_\infty-B_e)-\frac{E_e}{3r}-D_e-\frac e3A_e\right].
\end{equation}
For \(N=\norm{F}{6}\) and \(e\geq1\), elementary integration gives
\[
 |A_e|\leq CNe^{-2},\quad
 |B_\infty-B_e|\leq CNe^{-3/2},\quad
 |D_e|\leq CNe^{-1},\quad |E_e|\leq CN.
\]
In \(C_e\), if \(s\leq e/2\), then \(t\geq e/2\), and
\[
 \int_{e-s}^e\sqrt{s+t-e}\dd t=\frac23s^{3/2}.
\]
The two regions with one small input therefore contribute at most
\(CN^2e^{-7/2}\) to \(|C_e|/r\). With both inputs above \(e/2\), their
product is \(O(N^2e^{-6})\), and the integral divided by \(r\) contributes
\(O(N^2e^{-4})\). Hence \eqref{eq:remainder} is \(O(N^2r^{-7})\).
On bounded \(r\), \eqref{eq:origin} and the tail integrals give a
uniform \(CN^2\) bound. These estimates also establish the absolute
pointwise collision bounds used above. The mass correction contributes
only a Gaussian times a bounded quadratic scalar.

For radial profiles \(F_1,F_2\in X_6\), the symmetric polarization is
\[
 B_R(F_1,F_2)=\frac{R(F_1+F_2)-R(F_1-F_2)}4.
\]
Apply the diagonal bound to \(\lambda F_1\pm\lambda^{-1}F_2\), with
\(\lambda^2=\norm{F_2}{6}/\norm{F_1}{6}\) when both norms are nonzero.
It gives \(\norm{B_R(F_1,F_2)}{7}\leq
2C\norm{F_1}{6}\norm{F_2}{6}\); zero arguments are immediate.
Now \(R(F)-R(\widetilde F)=B_R(F-\widetilde F,F+\widetilde F)\)
proves \eqref{eq:R} and the smooth polynomial mapping properties.

For mass cancellation only the unweighted collision integral is needed.
Write \(V=(v+v_*)/2\), \(v-v_*=dk\), \(k\in\Sph^2\). The joint measure
including the collision factor is
\(\dd V\,d^3\dd d\,\dd k\,\dd\sigma\). Swapping \(k,\sigma\)
interchanges the pre- and post-collision velocities and preserves this
measure. The absolute loss integral is bounded by
\[
 4\pi\iint (|v|+|v_*|)|F(v)F(v_*)|\dd v\dd v_*<\infty;
\]
the gain has the same bound by this angular swap. Fubini now proves
\(m(Q(F,F))=0\). Together with \(m(CF)=0\), this proves \(m(R(F))=0\).
No energy-weighted collision interchange has been used.
\end{proof}

\section{The linear evolution and a flat tail}\label{sec:linear}

\subsection*{The evolution on the faster-decaying space}
We construct the correction evolution in \(X_p^0\) by limits of
compactly supported profiles. The critical \(r^{-6}\) tail is absent
from this space and will be prescribed separately.

\begin{lemma}\label{lem:semigroup}
There is a strongly continuous semigroup \(S(t)\) on \(X_p^0\) preserving
mass, with \(\norm{S(t)}{X_p^0\to X_p^0}\leq e^{bt}\) for a fixed \(b\).
If \(L\) is its generator, then
\(Lu(r)=\calL u(r)\) for \(u\in D(L)\), and \(M\in D(L)\).
\end{lemma}
\begin{proof}
Start with \(4\pi C_{\rm raw}\). On profiles supported in \([0,R]\),
its gain and loss are bounded operators, and the positive
variation-of-constants series defines the evolution. The gain uses only
larger radii, so zero extension is consistent as \(R\) increases. A
profile vanishing at \(R\) continues to vanish there.
Recall that \(p=13/2\). Where \(\chi>0\), \(r>4\), and
\begin{equation}\label{eq:row}
 \frac{(4/r)\int_r^\infty s g_p(s)\dd s}{r g_p(r)}
 =\frac4{p-2}(1+r^{-2})\leq\frac{17}{18}<1.
\end{equation}
At \(\chi=0\) both gain and loss terms vanish. Thus \(Ag_p\), for
\(A=\norm{F}{p}\), is a stationary supersolution bounding both signs of
the supported solution. Positivity of the bounded integral series
proves contractivity uniformly in \(R\). Density gives a positive
contraction semigroup \(U(t)\) on \(X_p^0\). Strong continuity and the
semigroup identity follow first on supported profiles and then by density.

The perturbation \(B=-4\pi\ell(\cdot)M\) is bounded on \(X_p^0\).
Set
\[
 S_0(t)=U(t),\qquad
 S_{j+1}(t)F=\int_0^t U(t-s)B S_j(s)F\dd s.
\]
These are vector-valued integrals for each fixed \(F\), with continuous
integrands. Their operator norms are bounded by
\(\norm{B}{}^jt^j/j!\). The sum is strongly continuous and satisfies
\(\norm{S(t)}{}\leq e^{\norm{B}{}t}\).
Uniqueness for the integral equation follows by iterating its norm
inequality; it also gives the semigroup identity. This construction
does not require operator-norm measurability of \(U(t)\).

For supported data,
\[
 m(U(t)F)=m(F)+4\pi\int_0^t\ell(U(s)F)\dd s.
\]
Boundedness of \(m,\ell\) extends this identity by density. In the
perturbed integral equation, the scalar mass derivative is
\(4\pi\ell(S(t)F)+m(BS(t)F)=0\). Thus \(S(t)\) preserves mass;
positivity of \(S(t)\) is not asserted.

At each fixed \(r\), evaluation of \(C_{\rm raw}F(r)\) is bounded in
the \(X_p\) norm. Passing from supported profiles to the integrated
equation gives the pointwise integrated identity for \(S\).
Taking the generator quotient gives \(Lu(r)=\calL u(r)\) on \(D(L)\).
Finally, smooth compact cutoffs \(M_j\) of the Gaussian satisfy
\(M_j\to M\) and \(C_{\rm raw}M_j\to C_{\rm raw}M\) in \(X_p^0\).
Pass to the limit in the integrated raw equation and divide by time;
the bounded perturbation then yields \(M\in D(L)\).
\end{proof}

\subsection*{The prescribed tail}
Let \(c=4\pi\), choose \(0<\varepsilon\leq1\), and define the flat function
\[
 a(t)=\begin{cases}\varepsilon e^{-1/t^2},&t>0,\\0,&t\leq0.\end{cases}
\]
For \(r>0\) and \(t\geq0\), set
\begin{equation}\label{eq:entrance}
 h(t,r)=\frac{c^5}{6r}\int_0^t a(s)(t-s)^4e^{-cr(t-s)}\dd s,\quad
 \widetilde H=\chi h,\quad
 \mu=m(\widetilde H),\quad H=\widetilde H-\mu M.
\end{equation}
Thus \(m(H)=0\) and \(\widetilde H\geq0\).

\begin{lemma}\label{lem:tail}
The map \(H\) is \(C^\infty\) in time into \(X_6\), with all time jets
zero at \(t=0\). For sufficiently small bounded times,
\begin{equation}\label{eq:Hbounds}
 \norm{H(t)}{6}\leq Ca(t),\qquad
 0\leq\mu(t)\leq Ct^3a(t),\qquad
 \lim_{r\to\infty}r^6H(t,r)=4a(t).
\end{equation}
The defect
\begin{equation}\label{eq:J}
 J=H_t-\calL H
\end{equation}
is \(C^\infty\) in time into \(X_p^0\), has zero mass and flat initial
jets, and satisfies \(\norm{J(t)}{p}\leq Ca(t)\).
Here \(\calL H\) is the pointwise expression; no assertion that
\(H\in X_p^0\) or \(H\in D(L)\) is made.
\end{lemma}
\begin{proof}
For \(k(u,r)=c^5u^4e^{-cru}/(6r)\),
\[
 \int_r^\infty s k(u,s)\dd s=\frac{c^4}{6}u^3e^{-cru},
 \qquad \partial_u k=cC_0k.
\]
Integration by parts in \(h=\int_0^\infty a(t-u)k(u,r)\dd u\)
therefore gives \(h_t=cC_0h\); the boundary at \(u=0\) vanishes.
After \(z=cru\),
\begin{equation}\label{eq:tail-integral}
 r^6h(t,r)=\frac16\int_0^\infty z^4e^{-z}a(t-z/(cr))\dd z
       \longrightarrow4a(t),
\end{equation}
by dominated convergence and \(\int_0^\infty z^4e^{-z}\dd z=24\).

For \(r\geq4\), the norm of this positive kernel as a map from bounded
scalar time functions to \(X_6\) is at most \(4(17/16)^3\).
Apply that bound to the uniform scalar difference-quotient remainders
of \(a\). It proves every \(X_6\)-valued time derivative and its flat
initial trace. This uses the supremum of the integrated kernel, not
the generally nonintegrable integral of its supremum.
Since \(a\) is increasing, it also gives the first bound in
\eqref{eq:Hbounds}. Integrating the positive kernel for mass and using
\(\int_0^\infty r e^{-cru}\dd r=(cu)^{-2}\) gives the bound on \(\mu\).
The Gaussian correction has no effect on the limit.

The expression
\[
 \widetilde H_t-cC_{\rm raw}\widetilde H
       =c\chi[C_0h-C_0(\chi h)]
\]
is supported in \(4\leq r\leq8\). All other terms in \(J\) are scalar
multiples of \(M\) or \(C_{\rm raw}M\), both rapidly decreasing.
Their coefficients are smooth by the kernel bound and bounded moment
functionals. For the amplitude bound, use \(h_t=cC_0h\) to obtain
\[
 |\mu'(t)|\leq C\int_4^\infty r^3h(t,r)\dd r\leq Ca(t).
\]
The same bound controls the compactly supported commutator term, \(\mu\), and
\(\ell(H)\). In particular no bound on \(a'/a\) is needed.
All these expressions are flat, lie in \(X_p^0\), and have zero mass
because \(m(H)=m(\calL H)=0\).
\end{proof}

\section{A smooth signed solution}\label{sec:signed}

Seek a mass-one profile in the form
\[
 F=M+H+w,\qquad w(0)=0,\qquad m(w)=0,\qquad w(t)\in X_p^0.
\]
Since \(Q(M,M)=0\), Lemma~\ref{lem:remainder} gives
\(R(M)=-\calL M=-LM\). Consequently the equation for the correction is
\begin{equation}\label{eq:w}
 w_t=Lw+\Phi(t,w),\qquad
 \Phi(t,z)=R(M+H(t)+z)-R(M)-J(t).
\end{equation}
Its mild form is
\begin{equation}\label{eq:mild}
 w(t)=\int_0^tS(t-s)\Phi(s,w(s))\dd s.
\end{equation}
All constants below are uniform for \(0<\varepsilon\leq1\) on a fixed
sufficiently short interval; they may depend on \(p,\chi,M\).

\begin{proposition}\label{prop:signed}
There is an interval \([0,\tau]\) on which \eqref{eq:mild} has a
mass-zero solution, with
\begin{equation}\label{eq:wbound}
 \norm{w(t)}{p}\leq C\int_0^t a(s)\dd s\leq Ct a(t).
\end{equation}
For every \(t_*<\tau\), the solution is \(C^\infty\) in time into
\(X_p^0\) on \([0,t_*]\), all initial time jets vanish, and
\(w(t)\in D(L)\). The profile \(F=M+H+w\) has mass one and satisfies
\(F_t=Q(F,F)\) pointwise.
\end{proposition}
\begin{proof}
In a fixed closed ball of \(C([0,\tau];X_p^0)\), the right side of
\eqref{eq:mild} has Lipschitz constant at most \(e^{b\tau}C\tau\).
For small \(\tau\) it is a contraction mapping the ball into itself.
The mass-zero constraint defines a closed subspace and is preserved by
both the semigroup and the forcing. The Banach fixed point theorem
therefore applies. The same estimates give
\[
 \norm{w(t)}{p}\leq C\int_0^ta(s)\dd s
                     +C\int_0^t\norm{w(s)}{p}\dd s.
\]
Gronwall, or iteration of this inequality and the factorial remainder,
proves \eqref{eq:wbound}.

We give the time regularity argument because the generator is
unbounded. The map \(\Phi\) is smooth into \(X_p^0\), quadratic in \(z\),
and every pure time derivative of \(\Phi(t,0)\) is zero at zero.
For \(h>0\), the semigroup property gives, for \(0\leq t\leq\tau-h\),
\begin{equation}\label{eq:difference}
 \frac{w(t+h)-w(t)}h
 =S(t)\frac{w(h)}h+
 \int_0^tS(t-s)
 \frac{\Phi(s+h,w(s+h))-\Phi(s,w(s))}h\dd s.
\end{equation}
The first term tends uniformly to zero by \eqref{eq:wbound}.
Taylor's formula for \(\Phi\) and the Lipschitz integral estimate
first bound these quotients and then show that they converge uniformly
on each strictly smaller time interval to the unique solution
\[
 y(t)=\int_0^tS(t-s)
 [\partial_t\Phi(s,w(s))+D_z\Phi(s,w(s))y(s)]\dd s.
\]
For clarity, the Taylor remainder tends uniformly to zero because the
quotients are uniformly bounded and the derivatives of the polynomial
map \(\Phi\) are uniformly continuous on the range under consideration.
Subtracting the equation for \(y\) from \eqref{eq:difference} and
iterating the integral inequality proves convergence without extraction.
Left derivatives agree at positive times. Hence \(w\) is \(C^1\),
\(w'(0)=0\), and \(w'\) satisfies the differentiated mild equation.

Repeat the argument for the differentiated system. At each order the
highest derivative enters linearly through \(D_z\Phi\); the other terms
use previously obtained derivatives. Every initial compatibility value
is zero because all earlier jets and all pure time jets of
\(\Phi(t,0)\) vanish. This proves all time derivatives, on the same
strictly interior interval, with flat initial traces.

Finally \eqref{eq:mild} implies
\[
 \frac{S(h)w(t)-w(t)}h
 =\frac{w(t+h)-w(t)}h
  -\frac1h\int_t^{t+h}S(t+h-s)\Phi(s,w(s))\dd s.
\]
The right side tends in \(X_p^0\) to \(w'(t)-\Phi(t,w(t))\).
This is exactly the generator-domain condition. Every forward quotient
is taken below \(\tau\); in particular it is valid at the final
displayed endpoint \(t_*<\tau\).
Using \(Lw=\calL w\), the definition of \(J\), and
\eqref{eq:split} gives \(F_t=Q(F,F)\) at each finite \(r\), including
zero by the continuous formulas. The mass is one by construction.
\end{proof}

\begin{proposition}\label{prop:smooth}
For every \(t_*<\tau\), the profile \(F(t,v)\) is jointly
\(C^\infty\) on \([0,t_*]\times\R^3\), including the velocity origin.
\end{proposition}
\begin{proof}
We already have \(F\in C^\infty_tX_6\). Recall the energy coordinate
\(e=|v|^2\) and its profile \(G(t,e)=F(t,\sqrt e)\). Therefore
\[
 A(t)=\int_0^\infty G(t,e)\dd e,\qquad
 D(t)=\int_0^\infty eG(t,e)\dd e
\]
are smooth scalar functions, since the two functionals are bounded on
\(X_6\). This argument uses no higher energy moments.
Fix \(E<\infty\). In \eqref{eq:gain}--\eqref{eq:loss}, replace
\[
 A_e=A(t)-\int_0^eG(t,s)\dd s,\qquad
 D_e=D(t)-\int_0^e sG(t,s)\dd s
\]
and use \eqref{eq:origin}. The equation on \([0,E]\) is now an ODE
on \(C([0,E])\) with the known smooth coefficients \(A,D\).
Its vector field \(\mathcal V_E(t,G)\) is locally Lipschitz on every
Banach space \(C^k([0,E])\) with its usual norm.
The kernels in \eqref{eq:origin} have fixed compact integration domains
and are integrable; differentiating \(G(eb)\) only introduces factors
\(b\in[0,1]\). The same holds for the other local integrals.

For \(k\geq1\), differentiation also gives the tame bound
\begin{equation}\label{eq:tame}
 \norm{\mathcal V_E(t,G)}{C^k}
 \leq C_{k,E}\bigl(|A|,|D|,\norm{G}{C^{k-1}}\bigr)
                     (1+\norm{G}{C^k})
\end{equation}
on bounded lower-norm sets. A derivative of order \(k\) in any product
is multiplied by an undifferentiated bounded factor; all other terms
involve orders at most \(k-1\). Polynomial factors in \(e\) are bounded
on the fixed interval.

Starting from \(G(0,e)=\pi^{-3/2}e^{-e}\), the local contraction in
\(C^k\) constructs a \(C^k\)-valued solution. Uniqueness in \(C^0\)
identifies it with the profile already obtained. Inductively,
\eqref{eq:tame} and Gronwall bound its \(C^k\) norm throughout the
existing time interval. This prevents earlier termination: the vector
field is bounded in \(C^k\), so a solution approaching a putative
termination time has a \(C^k\) limit and the local contraction restarts
it. Thus every \(k\) holds on the same interval.
The ODE, with its smooth coefficients, then gives all mixed time and
energy derivatives.

Finally \(F(t,v)=G(t,|v|^2)\) is smooth at \(v=0\) as well.
For a one-sided \(C^1\) energy profile,
\(\partial_{v_i}G(|v|^2)=2v_iG'(|v|^2)\), including at zero since
\([G(h^2)-G(0)]/h=O(h)\). Apply this identity successively to all energy
derivatives and use the polynomial product rule. No extension to
negative energies is required.
\end{proof}

\section{Positivity and the rapid-decay consequence}\label{sec:positive}

Positivity of the semigroup \(S(t)\) was not established in
Lemma~\ref{lem:semigroup}. We prove nonnegativity for the nonlinear
solution itself.
Shorten \(t_*<\tau\) whenever required below. With \(\mu\leq1/2\), the
baseline
\[
 P=M+H=(1-\mu)M+\widetilde H
\]
is nonnegative. For \(n=F_-=\max(-F,0)\) and
\(N(t)=\norm{n(t)}{p}\), \eqref{eq:wbound} gives
\begin{equation}\label{eq:n}
 0\leq n\leq|w|,\qquad N(t)\leq Ct a(t).
\end{equation}

\subsection*{Finite support of the negative part}
There are fixed \(K_0,k_0>0\) such that
\begin{equation}\label{eq:lower-tail}
 \widetilde H(t,r)\geq k_0a(t)r^{-6}
       \quad\text{when }r\geq K_0t^{-3}.
\end{equation}
To see this, restrict \(u=t-s\) in \eqref{eq:entrance} to
\([1/(cr),2/(cr)]\). For \(u\leq t/2\),
\[
 \log\frac{a(t-u)}{a(t)}
   =-\frac{2tu-u^2}{t^2(t-u)^2}\geq-\frac{8u}{t^3}.
\]
For \(r\geq K_0t^{-3}\) the restricted interval lies below \(t/2\),
and the cutoff equals one after a fixed shortening of time.
The change \(z=cru\) bounds \(r^6h/a(t)\) below by
\[
 \frac{e^{-16/(cK_0)}}6\int_1^2 z^4e^{-z}\dd z>0.
\]
On the same region the ratio of \(|w|\) to the lower bound
\eqref{eq:lower-tail} is at most
\[
 Ct r^{6-p}\leq C K_0^{6-p}t^{1+3(p-6)}\longrightarrow0.
\]
Hence \(F>0\) for \(r\geq R(t):=K_0t^{-3}\).
The negative part has compact support at every positive time, with one
common compact support on any compact positive-time interval.
Its weighted norm is continuous there, and tends to zero at zero by
\eqref{eq:n}. The full profile \(F\) is not claimed to have finite
\(X_p\) norm.

\subsection*{A strict mixed gain estimate}
Recall the Gaussian loss frequency
\[
 \nu_M(v)=4\pi\int_{\R^3}|v-w|M(w)\dd w.
\]
It is radial, so we also write \(\nu_M(r)\) for \(|v|=r\).
For every fixed \(\alpha\) with \(4/(p-2)<\alpha<1\), there is a
constant \(C_M\) such that every nonnegative radial \(n\leq Ng_p\)
satisfies
\begin{equation}\label{eq:mixed}
 \la r^p\,2Q^+(M,n)(r)
       \leq \alpha\nu_M(r)N+C_MN.
\end{equation}
To verify it, polarize \eqref{eq:gain}. The local mass integral of
the Maxwellian and the tail integral of \(n\) are
\[
 B_M(r)=\int_0^{r^2}\sqrt s\,G_M(s)\dd s,\qquad
 A_n(r)=\int_{r^2}^\infty G_n(s)\dd s.
\]
The principal term is \(8\pi^2(2/r)B_M(r)A_n(r)\), and
\[
 A_n(r)\leq\frac{2N}{p-2}(1+r^2)^{1-p/2}.
\]
After weighting and division by \(\nu_M(r)N\), its upper limit at
infinity is \(4/(p-2)\), since
\(\nu_M(r)/(8\pi^2r)\to B_M(\infty)\).
Every term containing the Maxwellian tail \(A_M\) decays exponentially
after polynomial weighting. In the mixed \(C\) term, split the
Maxwellian energy at \(e/2\). Below \(e/2\), the other input is at least
\(e/2\), giving, before division by \(r\), a bound
\[
 CN(1+e)^{-p/2}\int_0^\infty s^{3/2}G_M(s)\dd s.
\]
The weighted term is therefore \(O(N/r)\). Above \(e/2\), the
Maxwellian is exponentially small and all remaining integrals are
polynomially bounded. On bounded radii the formulas at zero give a
uniform \(C_MN\) bound. This proves \eqref{eq:mixed}.

We also need the elementary envelope bound
\begin{equation}\label{eq:bounded-gain}
 \norm{Q^+(B,n)}{\infty}\leq C\norm{B}{6}N
       \quad(B\geq0,\ 0\leq n\leq Ng_p).
\end{equation}
Indeed \(g_p\leq g_6\), so positivity of the gain bounds its left
side pointwise by
\(\norm{B}{6}NQ^+(g_6,g_6)\).
The three terms of \eqref{eq:gain} for this envelope are respectively
\(O(e^{-4}),O(e^{-5/2}),O(e^{-7/2})\) for \(e\geq1\); at zero use
\eqref{eq:origin}. Thus the envelope gain is bounded.

\subsection*{The compact maximum argument}
At a point where \(F<0\), expand \(F=F_+-n\) in the symmetric gain.
The nonnegative diagonal terms can be discarded to give
\begin{equation}\label{eq:negative-equation}
 n_t\leq2Q^+(F_+,n)-\nu_F n.
\end{equation}
Moreover \(F_+\leq M+\widetilde H+|w|\), and
\[
 \nu_F\geq(1-\mu)\nu_M-\nu_{|w|}
          \geq(1-\delta)\nu_M
\]
for any fixed small \(\delta>0\), after shortening time.
Here \(\nu_M\geq c_0\la v\), whereas
\(\nu_{|w|}\leq C\norm{w}{p}\la v\).
Fix \(\alpha=17/18\), \(\delta=1/36\). Since \(p=13/2\),
\(4/(p-2)=8/9<\alpha\) and \(1-\delta-\alpha=1/36>0\).

When \(N(t)>0\), its maximum is attained at a negative point with
\(r\leq R(t)\). Apply \eqref{eq:mixed} to the Maxwellian contribution
in \eqref{eq:negative-equation}. For the other gain terms,
\eqref{eq:bounded-gain} and
\(\norm{\widetilde H+|w|}{6}\leq Ca(t)\) give, after weighting on
this compact interval, at most \(Ca(t)t^{-3p}N(t)\).
Discarding the strictly negative remaining frequency term yields
\begin{equation}\label{eq:Dini}
 D^+N(t)\leq [C_M+Ca(t)t^{-3p}]N(t),
 \qquad D^+N(t)=\limsup_{h\downarrow0}\frac{N(t+h)-N(t)}h.
\end{equation}
This use of the derivative of a maximum is local: on a compact
positive-time interval all maximizing points lie in a common compact
radius interval. A maximizing sequence at \(t+h\) has a subsequence
converging to a maximizer at \(t\), and uniform integration of the
continuous time derivative bounds the upper Dini derivative by the
largest derivative at such maximizers.
If \(N(t)=0\), then \(F\geq0\); at every zero of \(F\), its equation
gives \(F_t=Q^+(F,F)\geq0\). The right derivative of the negative part
there is zero, as it is at strictly positive points. Thus
\eqref{eq:Dini} holds also in the zero-norm case.

The coefficient in \eqref{eq:Dini} is integrable at zero, since the
flat function \(a(t)t^{-3p}\) tends to zero faster than any fixed
power. For \(0<\eta<t\), the scalar Dini comparison gives
\[
 N(t)\leq N(\eta)
 \exp\left(\int_\eta^t[C_M+Ca(s)s^{-3p}]\dd s\right).
\]
For completeness, this comparison follows by multiplying \(N\) by
the integrating factor. A continuous function with nonpositive upper
right Dini derivative is nonincreasing: otherwise subtract a small
positive linear function and take a minimum before an interval's right
endpoint to contradict the strictly negative Dini derivative.
Now \eqref{eq:n} implies \(N(\eta)\to0\), while the exponential remains
bounded. Thus \(N=0\) and \(F\geq0\).

\subsection*{The macroscopic bounds and the endpoint}
All the preceding reductions select a single positive
\(t_*<\tau\). Set \(f(t,x,v)=F(t,|v|)\).
Propositions~\ref{prop:signed} and \ref{prop:smooth}, and positivity,
give the claimed smooth nonnegative solution on this interval.
Its datum is \(M\). Every velocity derivative of \(M\) is a polynomial
times that Gaussian, and every spatial derivative is zero.
The time jets of \(H+w\) vanish at zero as well.

There is a fixed \(K\geq1\) with
\[
 0\leq F(t,r)\leq K(1+r^2)^{-3}\qquad(0\leq t\leq t_*).
\]
The mass is one, and the energy is at most
\[
 E_0=4\pi K\int_0^\infty r^4(1+r^2)^{-3}\dd r<\infty.
\]
The positive entropy is bounded by \(\log K\int F=\log K\).
To check that the entropy integral is well defined, for \(b>0,z\geq0\)
use
\[
 -z\log z\leq b(1+|v|^2)z+e^{-1-b(1+|v|^2)}.
\]
The scalar inequality follows by maximizing \(-z\log z-cz\) with
\(c=b(1+|v|^2)\). Its Gaussian remainder is integrable, so the negative
entropy part is uniformly finite as well. No entropy dissipation law
is required.

Finally the Gaussian and \(w\in X_p^0\), \(p>6\), have zero \(r^6\)
tail. Equations \eqref{eq:Hbounds} and \eqref{eq:tail-integral} give
\eqref{eq:main-tail}. This proves Theorem~\ref{thm:construction}.

\begin{corollary}\label{cor:literal}
For the full-sphere hard-sphere equation \eqref{eq:boltzmann} on
\(\T^3\times\R^3\), joint pointwise smoothness, initially Schwartz
data with Schwartz derivatives, and the uniform local bounds
\eqref{eq:macroscopic} do not imply an extension with rapid velocity
decay at the endpoint. More precisely, for any \(0<T<t_*\), the
restriction of the solution in Theorem~\ref{thm:construction} to
\([0,T)\) satisfies these hypotheses with \(m_0=M_0=1\) and finite
\(E_0,H_0\), but every continuous extension to \(T\) has
\[
 \lim_{r\to\infty}r^6f(T,x,r\omega)
       =4\varepsilon e^{-1/T^2}>0.
\]
\end{corollary}
\begin{proof}
At each finite velocity an extension must have the existing smooth
endpoint value, so its value at \(T\) is precisely the constructed
profile. That profile has the displayed tail and cannot be rapidly
decreasing.
\end{proof}

Here \(t_*,\varepsilon\) are those of Theorem~\ref{thm:construction},
\(T\) is any strictly interior positive time, \(r\) is velocity radius,
\(\omega\in\Sph^2\) is a unit direction, and \(x\in\T^3\).
The solution already extends smoothly in \(X_6\) across \(T\).
Thus the failed conclusion is rapid velocity decay. The example does
not resolve the continuation question in a conservative,
energy-admissible, or positive-time Schwartz solution class.

\paragraph{AI disclosure.}
OpenAI Codex was used for proof development, calculation checks, and
drafting. The mathematical claims and the limitations of their scope
are stated in the text. One fresh-context AI reader and one fresh-context
AI referee provided feedback on a preceding draft; the Codex writing
assistant then revised it once.

\begin{thebibliography}{9}
\raggedright
\bibitem{silvestre}
L. Silvestre, \emph{Regularity estimates and open problems in kinetic
equations}, arXiv:2204.06401v1 (2022).
\href{https://arxiv.org/abs/2204.06401v1}{arXiv:2204.06401v1}.

\bibitem{lu-wennberg}
X. Lu and B. Wennberg, \emph{Solutions with increasing energy for the
spatially homogeneous Boltzmann equation},
Nonlinear Analysis: Real World Applications \textbf{3} (2002), no.~2,
243--258.
\href{https://doi.org/10.1016/S1468-1218(01)00026-8}
{doi:10.1016/S1468-1218(01)00026-8}.

\bibitem{basile}
G. Basile, D. Benedetto, L. Bertini, and E. Caglioti,
\emph{Asymptotic probability of energy increasing solutions to the
homogeneous Boltzmann equation}, arXiv:2202.07311v2 (2022).
\href{https://arxiv.org/abs/2202.07311v2}{arXiv:2202.07311v2}.
\end{thebibliography}
\end{document}
